% !TeX root = ../教材题库_全部条目_转写工作稿.tex
% Source: XJ-B1-C06-HW-P255-Y1208
% 第6章 / 6.4 用样本估计总体 / 教材 p.255 / PDF 第262页
% OCR未校对：公式与文字均可能有误；未完成数学TeX转写；禁止直接发给学生
% 原题及相邻原解截图（核对用）：../题目原图/XJ-B1-C06-HW-P255-Y1208-01.webp
% 本文件只在转写完成后作为题干使用；OCR 条目尚不是可靠题干。
\begin{bankitem}{习题6.4 第7题}{第6章 / 6.4 用样本估计总体 / 教材 p.255 / PDF 第262页}{XJ-B1-C06-HW-P255-Y1208}{OCR未校对：公式与文字均可能有误；未完成数学TeX转写；禁止直接发给学生}
\begin{ocrdraft}
[教材p.255 / PDF第262页]
\par
7．研究人员随机调查了某地1000
\par
组距
\par
名“上班族”每天在工作之余使用手机
\par
0，20
\par
上网的时间，并绘制了如图所示的频率
\par
分布直方图．若同一组中的数据用该组
\par
0．12
\par
0．10
\par
区间的组中值作代表，试估计该地“上
\par
0．08
\par
班族”每天在工作之余使用手机上网的
\par
平均时间．
\par
4
\par
6
\par
8
\par
时间/h
\end{ocrdraft}
\end{bankitem}
